\section{Context 管理：最难的问题}

\subsection{为什么 Context 管理如此困难？}

文章作者坦言：\textbf{最难的问题不是调用 LLM，而是管理送入 LLM 的内容}。

Utah 使用的 Tool 可能每次调用返回数千字符。经过几轮迭代后，对话上下文急剧膨胀，模型开始"迷失"——作者观察到 Agent 陷入无限循环地调用 Tool 却永远不产出最终回复。

\begin{warningbox}{Context 膨胀的恶性循环}
Tool 返回大量内容 $\to$ context window 被填满 $\to$ 模型无法关注关键信息 $\to$ 模型反复调用 Tool 试图"理解"情况 $\to$ 更多内容涌入 $\to$ 恶性循环。这不是 LLM 的 bug，而是 context window 有限容量下的必然结果。\textbf{任何使用 Tool 的 Agent 都必须面对这个问题}。
\end{warningbox}

\subsection{两层 Context 剪枝}

Utah 的解决方案是\textbf{两层 context pruning}：

\begin{lstlisting}[caption={Context 剪枝配置},language=Java]
const PRUNING = {
  keepLastAssistantTurns: 3,
  softTrim: {
    maxChars: 4000,
    headChars: 1500,
    tailChars: 1500
  },
  hardClear: {
    threshold: 50000,
    placeholder: "[Tool result cleared]"
  },
};
\end{lstlisting}

具体策略如下：

\begin{table}[H]
\centering
\caption{Context 管理策略}
\begin{tabular}{lp{9cm}}
\toprule
\textbf{策略} & \textbf{描述} \\
\midrule
Soft Trim & 旧的 Tool 结果保留头部 1500 字符 + 尾部 1500 字符，中间截断 \\
Hard Clear & 当总 context 超过 50,000 字符时，旧的 Tool 结果完全替换为占位符 \\
最近保留 & 最近 3 轮 assistant 的输出始终完整保留 \\
\bottomrule
\end{tabular}
\end{table}

\subsection{跨运行的 Session Compaction}

Pruning 处理的是\textbf{单次运行内}的 context 管理。但 Agent 可能跨多次运行与用户持续对话。Session compaction 处理的是\textbf{跨运行}的上下文积累：当估算的 token 数超过阈值时，对话历史会被\textbf{总结压缩}后再送入下一次运行。

此外，Utah 还实现了两个辅助机制：

\begin{itemize}
    \item \textbf{Budget Warning}：当 Agent 接近最大迭代次数时，注入系统消息告诉它"请尽快收尾"
    \item \textbf{Overflow Recovery}：如果 LLM 在运行中返回 context-too-large 错误，强制执行 compaction 并重试，不浪费迭代次数
\end{itemize}

\begin{knowledgebox}{Context 管理的四重防线}
Utah 的 context 管理可以总结为四重防线：
\begin{enumerate}
    \item \textbf{Pruning}（单次运行内）：裁剪旧的 Tool 结果
    \item \textbf{Compaction}（跨运行）：总结压缩对话历史
    \item \textbf{Budget Pressure}（预防性）：提前警告 Agent 收尾
    \item \textbf{Overflow Recovery}（应急性）：context 超限时强制压缩并重试
\end{enumerate}
这四重防线共同确保 Agent 不会因 context 膨胀而"迷失"。
\end{knowledgebox}

\subsection{本章小结}

Context 管理是 Agent 工程中最被低估的挑战。它不是一个可以一次性解决的问题，而是需要多层防线协同工作的系统性方案。Utah 的实践表明，pruning、compaction、budget warning 和 overflow recovery 四管齐下，才能让 Agent 在长时间运行中保持稳定。


